\section{Fourteen unrestricted integer unknowns}
\label{sec:signed14}

We combine package~30's quartic in $31$ positive unknowns, package~27's
fast Jones evaluation, and package~29's signed square-stack and intersective
finalizer. Two new restricted-completeness arguments and a Pell orbit remove
three of package~29's guards. The resulting fixed polynomial has
\[
 \boxed{14\text{ integer unknowns},\quad
 \deg F_{14}=813056\cdot5^{32}+2752,\quad1524\text{ binary operations}.}
\]
A lower-degree alternative has $16$ integer unknowns, degree
$406528\cdot5^{32}+10008$, and $1537$ operations. These are ordinary,
unformalized reductions, not optimality claims or a proof of consistency.
The complete derivation and exact source boundaries are in
\texttt{signed/SIGNED\_PROOF.md}.

\subsection{Fixed index and the exact polynomial}
Use the \emph{whole} compact943 coefficient tuple of package~30 and its
normalized quartic $P_{31}(x,\mathbf n)=\sum_jr_j(\mathbf n+1)^2$.
The dummy $x$ is fixed to $1$. Put $L=5^{32}$ and
\[
 c_\alpha=(4-|\alpha|)!\prod_j\alpha_j!\,[\mathbf X^\alpha]P_{31},
 \qquad z=2^{45933757},
\]
\[
 e(T)=\sum_{|\alpha|\le4}(z+c_\alpha)
 T^{L-\sum_{j=0}^{31}\alpha_j5^j},\qquad y=e(2z).
\]
The $402$ nonzero coefficient records are in
\texttt{signed/upstream/}\allowbreak\texttt{COMPACT\_CONSTANTS.json}. Unlisted coefficients are zero; the index
certificate proves $|c_\alpha|+1<z$.
Consequently $y<(2z)^{L+1}$. Define the additional fixed integer
\[
 \kappa=45933758(L+1)+1,\qquad a=2^\kappa.
\]
Thus $\kappa$ is odd and $a>y$. This enormous coefficient is given by a
bounded recipe, \emph{not materialized}. The same applies to $y$ and the
weighted coefficients. Fixed leaves are not existential variables.

The complete list of unknowns is
\[
 (\varepsilon,g,h,s,w,\phi,i,j,z_1,z_2,z_3,z_4,v,u)\in\mathbb Z^{14}.
\]
Set $b=a\varepsilon^2$, $\omega=w^2$ and define the following
\emph{polynomial abbreviations}, none of which is an extra unknown:
\[
 B=b^5,\quad q=B^L,\quad \ell=\sum_{r=1}^{31}B^{5^r},\quad
 \lambda=\sum_{r=0}^{4L-1}B^r,\quad \theta=B-2z,
 \quad N=q^{16},\quad G=32bB^{5^{31}}.
\]
\begin{align*}
 R={}&\left[g+e(B)q^3+\ell q^5+
 \left(2(e(B)-z\lambda)(1+B+g)^4+\lambda B+\lambda Bq^4\right)q^7\right](N^2-N)\\
 &+\left[q^3-b\ell+\ell+\theta\lambda q^3+(B-2)q^8\right](N^2-1).
\end{align*}
\begin{gather*}
 U=N^2\omega,\quad Y=N^2s,\quad M=RY,\quad P=2M^2U,
 \quad K=R+1+h(P-1),\quad A=M(U+1),\\
 B_1=2R+1,\quad W_1=b\omega,\quad C=B_1+W_1+\phi,
 \quad J=4A-5,\quad Q=3W_1C-2(W_1^2-1),\\
 D=(A^2-1)C^2+1,\quad E=iJDC^2,\quad F=(A^2-1)E^2+1,\\
 G_1=A+F(CD+1-A),\quad H=B_1+jC,\quad I=(G_1^2-1)H^2+1,\\
 a_0=(P^2-1)K^2+1,\quad a_1=DFI,\quad
 m=K^2G-4(C-KY)^2G-8gK^2,\\
 D_0=JF,\qquad T_0=QF+J(H-C).
\end{gather*}
Put $(d_1,d_2,d_3,d_4)=(3g-1,3s-1,3\phi+2,3m-1)$ and
\[
 C_0=1,\qquad B_j'=1+4d_jC_{j-1}^2z_j^2,\qquad
 C_j=C_{j-1}B_j'\ (1\le j\le4),\qquad C_*=C_4.
\]
Finally set
\begin{gather*}
 Z=\varepsilon wi z_1z_2z_3z_4,\quad S=ZD_0^5,\quad r=T_0+D_0v,\\
 T=(r^3-D_0^2r-D_0^3)(r^2+23D_0^2),\quad
 w_1=C_*+1,\quad w_2=w_1(P^2+K^2+2),\quad X=Su-T,\\
 A_0=S^2C_*,\quad A_1=S^2w_1^2a_0,\quad A_2=S^2w_2^2a_1,\\
 V=X^2+A_2-A_0-A_1,\qquad W=V^2+4X^2A_2-4A_0A_1.
\end{gather*}
The requested equation is the single integer polynomial equation
\[
 \boxed{F_{14}=W^2-16V^2X^2A_2=0.}
\]
Its division-free expanded arithmetic DAG is
\texttt{signed/generated/}\allowbreak\texttt{signed14.json}. The exported names
\texttt{residue} and \texttt{normroot} mean $v,u$.

\subsection{The signed Jones interface and a removable guard}
The imported Jones ratio theorem, on its original domain
$b=y+1+\varepsilon_0$, $\varepsilon_0,\phi\ge0$ and
$g,h,s,\omega,i,j\ge1$, identifies quartic solvability with
\[
 a_0,a_1\in\operatorname{Sq}(\mathbb Z),\qquad D_0\mid T_0,
 \qquad m>0.\tag{JR}
\]
For every integer assignment $D_0\ne0$: $J=4A-5\ne0$, and $F=0$
would force $A=0,E=\pm1$, although $E$ would be divisible by $5$.
Replacing a nonzero $i$ by $|i|$ leaves (JR) unchanged since only $E^2$
occurs in its consumers.

Package~29's proof of automatic $h$ positivity does not use $j$.
For completeness, assume $b>y$, $g,s,\omega\ge1$, $\phi\ge0$,
$a_0$ square and $m>0$. The margin gives
\[
 K\ne0,\quad g<4bB^{5^{31}},\quad 4(C-KY)^2<K^2.
\]
Writing $t=5^{31}$, one has $1+B+g<B^{t+2}\le q$ and
$b\ell<q$. In the first bracket for $R$, the potentially negative
term is absorbed by $\lambda Bq^4>2z\lambda(1+B+g)^4$.
Both brackets are positive, so $R\ge2$, $C>0$, and
$P\ge2R^2>R+2$. The margin implies $K>0$. If $h\le0$, then
$1\le K\le R+1<P$. The case $K=1$ would imply $P-1\mid R$;
for $2\le K<P$ the radicand lies strictly between $(PK-1)^2$ and
$(PK)^2$. Both contradict the hypotheses. Hence $h\ge1$.

\paragraph{New Pell-orbit lemma.}
Under those hypotheses and (JR), $A\ge4$, $D\ge A^2$, $F>0$ and
$G_1>A>1$. Since $F\equiv1\pmod D$ and $F\equiv1\pmod J$,
the divisibility gives $H\equiv C\pmod F$. Thus
\[
 G_1\equiv1\pmod D,\quad I\equiv1\pmod D,
 \qquad G_1\equiv A\pmod F,\quad I\equiv D\pmod F.
\]
Therefore $D,F,I$ are positive and pairwise coprime. As $DFI$ is a
square, all three are squares. Put $d=G_1^2-1$, a positive nonsquare,
and $x=\sqrt I>0$. Even if $H<0$, $x+H\sqrt d>0$.
The unit $\eta=G_1+\sqrt d>1$ has norm one, so some positive power
is the identity modulo $CF$ in the integer quadratic algebra.
Multiplying $x+H\sqrt d$ by arbitrarily large such powers yields
\[
 x_n^2-dH_n^2=1,\qquad H_n\equiv H\pmod{CF},\qquad H_n\longrightarrow+\infty.
\]
Choose $H_n>B_1$ and put $j_n=(H_n-B_1)/C>0$. This preserves
$a_1$ being square and $D_0\mid T_0$, and leaves $a_0,m$ unchanged.
Thus the positivity guard on $j$ can be deleted entirely.

\subsection{Cofinal bases and square witnesses: two more deleted guards}
The proof of \texttt{USys\_of\_mem} in pinned
\texttt{Master3.lean}, lines~303--510, chooses an exponent
$k=xy+\sum n_j+1$ only to obtain $k\ge1$, $xy<2^k$, and
$n_j<2^k$ for a fixed quartic witness. The rest of the construction
uses those inequalities, not that equality. Hence it works for every
sufficiently large exponent $k$, including any prescribed parity.
\texttt{NineReduce.reduced\_of\_mem} and
\texttt{NineRatio.ratio\_complete} transport that same base $b=2^k$.
At $x=1$ we may therefore choose $k=\kappa+2r$ large enough, and set
$\varepsilon=2^r$, giving $b=a\varepsilon^2$. Conversely any
nonzero integer $\varepsilon$ gives $b\ge a>y$, so the old ceiling
coordinate $\varepsilon_0=b-y-1$ is nonnegative. This is a cofinal
restriction of the base, not an unjustified square substitution for
$\varepsilon_0$.

There is a second concrete freedom. In pinned \texttt{Ratio.lean},
\texttt{exists\_B25core}, lines~434--442, the canonical witness is
\[
 \omega=2^{2R+1-k},\qquad U=N^2\omega.
\]
The source first proves $k<2R+1$ from $k<2^k=b\le R$, so natural
subtraction does not truncate. Since our $k$ is odd,
$2R+1-k$ is positive and even: this $\omega$ is a positive square.
\texttt{NineRatio.ratio\_complete}, lines~270--279, retains that
same $\omega$ while adding $i,j$ and changing $\phi$. Thus
$\omega=w^2$, $w\ne0$, preserves completeness; soundness follows
from $w^2\ge1$. There is no analogous asserted square property for $s$.

Consequently only $g,s\ge1$, $\phi\ge0$ and $m>0$ require inequality
guards. Nonvanishing of $\varepsilon,w,i$ is sufficient for the other
restrictions. These source-inspected strengthenings are ordinary proofs,
not newly compiled formal theorems.

\subsection{Why the finalizer has exactly the intended zeros}
Every $B_j'$ is $1$ modulo $4$ and coprime to $C_{j-1}$.
If all $z_j\ne0$ and $C_*$ is a square, unique factorization forces
$|B_j'|$ to be squares; their residue modulo $4$ excludes negative
squares. Thus $d_j\ge0$, and $d_j\equiv2\pmod3$ excludes zero.
These are exactly the four required inequalities. Conversely each
positive $d_j\equiv2\pmod3$ is nonsquare. A Pell unit power gives
$u_j^2-d_jv_j^2=1$ with $v_j$ divisible by $2|C_{j-1}|$.
Taking $z_j=v_j/(2C_{j-1})\ne0$ makes every $B_j'$ a square.

The polynomial $\Pi(t)=(t^3-t-1)(t^2+23)$ has no rational root but
has a root modulo every nonzero integer. To see the latter, for an
odd prime other than $23$, either the cubic has a simple root, or
its irreducibility makes Frobenius a three-cycle and its discriminant
$-23$ a nonzero square. Thus one factor has a simple root and lifts.
At $23$ use the cubic root $3$; at $2$ use $-23\equiv1\pmod8$ and
the standard square-root lifting. Apply the Chinese remainder theorem.
For $D\ne0$ the homogeneous expression
\[
 \mathcal R(D,T,v)=((T+Dv)^3-D^2(T+Dv)-D^3)
 ((T+Dv)^2+23D^2)
\]
is nonzero and satisfies
\[
 (D\mid T\ \&\ Z\ne0)\quad\Longleftrightarrow\quad
 \exists v:\ ZD^5\mid\mathcal R(D,T,v).
\]
The forward direction applies intersectivity after $T=Dq$.
For the converse, if $v_p(D)=\alpha>\beta=v_p(T)$, the two factors
have valuations $3\beta,2\beta$, contradicting divisibility by $D^5$.
The case $Z=0$ would contradict nonvanishing. This supplies exactly
$D_0\mid T_0$ and all nonzero conditions via $S\mid T$.

The new linear weight is valid even before recovering signs:
$C_*\equiv1\pmod4$ implies
\[
 |C_*+1|>\sqrt{|C_*|},\qquad
 \sqrt{|a_0|}\le P^2+K^2+1.
\]
For negative $C_*=-t$, necessarily $t\ge3$ and
$(t-1)^2-t=t^2-3t+1>0$. Hence
\[
 |w_2|>\sqrt{|C_*|}+|w_1|\sqrt{|a_0|}.
\]
For complex roots $r_0^2=C_*,r_1^2=a_0,r_2^2=a_1$, multiplication gives
\[
 F_{14}=\prod_{\sigma\in\{-1,1\}^3}
 \bigl[X-S(\sigma_0r_0+\sigma_1w_1r_1+\sigma_2w_2r_2)\bigr].
\]
If $S=0$, this is $T^8\ne0$. Otherwise a zero makes one encoded
root sum rational. The strict absolute-weight bounds separate all
sign choices on nonzero roots, including imaginary roots. Every
Galois automorphism of their multiquadratic field must therefore fix
each root, so all three roots are rational, hence integers. It follows
that $S\mid T$ and $C_*,a_0,a_1$ are integer squares. Conversely these
conditions yield a zero by taking
$u=T/S+r_0+w_1r_1+w_2r_2$. Zero roots and repeated square classes
cause no exception.

Combining this exact elimination with the square stack, homogeneous
merger, recovered $h$, sign change of $i$, and Pell orbit for $j$
recovers the original Jones witness. Conversely the cofinal-base and
square-$\omega$ constructions, then the stack and finalizer, produce
an integer zero from any quartic witness. The imported compiler and
literal K215 calibration give
\[
 \mathrm{ZFC}\vdash
 \bigl(\nexists\mathbf t\in\mathbb Z^{14}:F_{14}(\mathbf t)=0\bigr)
 \ \longleftrightarrow\ \operatorname{Con}(\mathrm{ZFC}).
\]

\subsection{Lower-degree alternative and exact complexity}
For $F_{16}$ keep $b=y+1+\varepsilon$, $\omega=w$ and use six
stack discriminants
\[
 3\varepsilon+2,\ 3g-1,\ 3s-1,\ 3w-1,\ 3\phi+2,\ 3m-1,
 \qquad Z=i\prod_{r=1}^6z_r.
\]
The same finalizer and signed linear weights apply. Its proof uses
only the new $j$-orbit and weight improvement, not either cofinal
specialization. The intermediate $15$-variable construction restricts
only $b$; its $1530$ operations and degree $813056L+4184$ are dominated
by $F_{14}$.

The fast evaluator shares fifth powers, evaluates the background digit
sum by the complete homogeneous polynomial $h_4$, and uses fixed
multivariate Horner evaluation for sparse corrections. In particular,
$1+t+t^2+t^3+t^4=1+(t^2+t)(t^2+1)$ saves a gate when $t^2$ is shared.
All scalar multiplications and squares are charged; fixed coefficient
construction is excluded, as in the comparison convention.
For $F_{14}$,
\[
\begin{array}{c|c@{\qquad}c|c}
 R&470L&P&1900L+4\\
 K&1900L+5&a_1&49280L+98\\
 m&4442L+14&C_*&4442L+133\\
 S&68750L+152&w_2&8242L+143\\
 A_2&203264L+688&F_{14}&813056L+2752
\end{array}
\]
Here $\deg A_2>\deg X^2,\deg A_0,\deg A_1$, so the successive
leading terms of $V,W,F_{14}$ are $A_2,A_2^2,A_2^4$.
A nonzero symbolic diagonal leading coefficient, certified modulo
$1000003$ as a scalar times a power of the positive fixed $a$, proves
these are exact degrees. For $F_{16}$ one obtains
$\deg A_2=101632L+2502$.

The delivered generator and checker reproduce the DAGs, counts, and
degrees. Diagnostics include $216$ actual-coefficient modular circuit
comparisons, $1000$ eight-factor norm evaluations, $840$ coprimality
cases, $120$ negative-to-positive Pell orbits, and intersectivity tests
for every modulus $1$ through $1200$. These checks supplement, and do
not replace, the universal proofs above.
